% TOP_PROBLEM: {"problemId": "problem.residual-finiteness-of-hyperbolic-groups", "canonicalTitle": "Residual Finiteness of Hyperbolic Groups", "releaseRank": 143, "primaryDomain": "algebra_representation_category", "primaryDomainLabel": "Algebra, representation & category theory", "displayStatus": "open_with_solved_subcases", "releaseStatus": "open_with_solved_subcases"}
% !TeX program = lualatex
% Definition and sources only; this file is not an LLM solution attempt.
\documentclass[11pt]{article}
\usepackage[margin=25mm]{geometry}
\usepackage{fontspec}
\setmainfont{Latin Modern Roman}
\usepackage{amsmath,amssymb,mathtools}
\usepackage{unicode-math}
\setmathfont{Latin Modern Math}
\usepackage{xurl,hyperref}
\hypersetup{colorlinks=true,urlcolor=blue}
\setcounter{secnumdepth}{0}
\setlength{\parindent}{0pt}
\setlength{\parskip}{5pt}
\setlength{\emergencystretch}{3em}
\begin{document}
\section*{\texorpdfstring{Residual Finiteness of Hyperbolic Groups}{Residual Finiteness of Hyperbolic Groups}}\label{143}
\subsection{Definitions and mathematical statement}
A group $G$ is residually finite if its finite quotients separate its elements from the identity:
\[
 \forall g\in G\setminus\{1\}\ \exists F\text{ finite}\ \exists\varphi:G\to F:
 \varphi(g)\ne1.
\]
The conjecture asks for this property for every word-hyperbolic group. A different quotient may be used for each $g$, and no algorithm or quantitative bound on $|F|$ is required. Hyperbolicity is defined using a finite generating set and the thin-triangle condition on its Cayley graph.
\par\smallskip\textit{Formulation and concept references:} \href{https://ems.press/content/serial-article-files/49988}{[S1]}.

\subsection{Short English statement}
Can every nonidentity element of every hyperbolic group be detected in some finite quotient?
\subsection{Sources}
\begin{itemize}
\item[S1]Nicolas Tholozan and Konstantinos Tsouvalas, “Residually Finite Non-Linear Hyperbolic Groups,” Commentarii Mathematici Helvetici 100 (2025), 1–9, DOI 10.4171/CMH/581.
\url{https://ems.press/content/serial-article-files/49988}.
\textit{Formulation source.}
\end{itemize}
\end{document}
